\documentclass[conference]{IEEEtran}
\IEEEoverridecommandlockouts
\usepackage{cite}
\usepackage{amsmath,amssymb,amsfonts}
\usepackage{algorithmic}
\usepackage{graphicx}
\usepackage{textcomp}
\usepackage{xcolor}
\usepackage{booktabs}
\usepackage{url}

\def\BibTeX{{\rm B\kern-.05em{\sc i\kern-.025em b}\kern-.08em
    T\kern-.1667em\lower.7ex\hbox{E}\kern-.125emX}}

\begin{document}

\title{Evaluating Seasonal and Disruption-Aware Ranking for Sustainable Nature and Heritage Tourism Recommendation\\
\thanks{Academic Minor Project, Department of Computer Science \& Engineering.}
}

\author{\IEEEauthorblockN{Student Researcher}
\IEEEauthorblockA{\textit{Department of Computer Science and Engineering} \\
\textit{BharatVerse Project Group}\\
India \\
email@institution.edu}
}

\maketitle

\begin{abstract}
This prototype compares popularity-based, seasonal, and disruption-aware ranking across nine Indian waterfall and trekking destinations. The benchmark contains 36 scenarios: 24 category-specific and 12 mixed-pool scenarios. Destination-month inputs and Grades 0--3 relevance labels are provisional; labels were assigned by one project annotator, independent review has not occurred, and complete record-level provenance is pending.

Overall macro-average NDCG@K is 0.6143 for Method A and 0.9776 for Methods B and C. Category-specific NDCG@3 is 0.6483, 0.9789, and 0.9789, respectively. No supplied disruption record passes the evidence policy, so closure-conflict rates are not estimable. Method C ties Method B because no disruption is applied. These descriptive results do not establish real-world safety, source validity, or generalization.
\end{abstract}

\begin{IEEEkeywords}
Tourism Recommendation Systems, Context-Aware Recommenders, Seasonal Travel Planning, Explainable AI, Safety Invariants, Ranking Evaluation, NDCG.
\end{IEEEkeywords}

\section{Introduction}
Tourism recommendation systems (TRSs) play an increasingly vital role in steering tourist dispersion, economic livelihood, and cultural engagement \cite{gavalas2014}. With the rapid growth of algorithmic travel aggregators, users increasingly trust automated engines to recommend sightseeing destinations, hiking trails, and heritage landmarks. However, modern commercial TRSs are almost ubiquitously optimized for engagement metrics, historical popularity, and collaborative filtering patterns \cite{werthner2004}.

While popularity-based heuristics succeed in static urban contexts (e.g., museums, city monuments), they introduce catastrophic hazards when applied to nature, adventure, and seasonal heritage tourism. In countries with extreme seasonal transitions such as India---governed by intense Southwest and Northeast monsoons, severe Himalayan winters, and blistering pre-monsoon heat waves---a destination's appeal and safety are intrinsically time-dependent. 

\subsection{Problem Statement and Popularity Inertia}
Consider two benchmark scenarios encoded in the supplied inputs (the disruption records have not been independently verified):
\begin{enumerate}
    \item \textbf{A supplied winter-closure scenario}: In the benchmark inputs, Valley of Flowers has a high popularity score and an absolute-closure record for November through May. The record is unverified in the repository; this illustrates algorithm behavior on an input record, not independently established access conditions.
    \item \textbf{A supplied monsoon restriction scenario}: The dataset marks Dudhsagar with an absolute-closure record in July and August. That record is marked unverified, so the benchmark does not establish the notice, its scope, or current access conditions.
\end{enumerate}

We define this systematic failure as \textbf{Popularity Inertia}: \textit{the tendency of review-weighted recommenders to push high-volume destinations into top recommendation slots irrespective of transient or seasonal unviability.}

\subsection{Research Questions}
To systematically investigate this phenomenon, we formulate three core Research Questions:
\begin{itemize}
    \item \textbf{RQ1 (Baseline Vulnerability)}: To what extent does static popularity-based ranking recommend unviable or legally closed destinations across seasonal variations in nature tourism?
    \item \textbf{RQ2 (Environmental Suitability)}: How do the methods compare in NDCG on the provisional benchmark labels?
    \item \textbf{RQ3 (Disruption-Aware Ranking)}: Does Method C exclude items that match absolute-closure records supplied to the benchmark, and what ranking metrics result?
\end{itemize}

\section{Related Work}
Early TRSs categorized points of interest using collaborative filtering (CF) and matrix factorization \cite{koren2009}. While CF captures latent user preferences, it notoriously suffers from cold-start problems and popularity bias, disproportionately favoring hyper-frequent items \cite{abdollahpouri2017}. In tourism, popularity bias is not merely a diversity issue; it is a direct driver of overtourism, environmental strain, and traveler dissatisfaction.

Context-Aware Recommender Systems (CARS) incorporate contextual variables such as time, location, companion type, and weather \cite{adomavicius2011}. Baltrunas et al. explored weather-aware mobile recommendations, demonstrating that rain significantly alters dining and indoor museum preferences \cite{baltrunas2012}. However, prior context-aware tourism models typically treat weather as a continuous soft preference penalty rather than treating official administrative warnings and legal closures as hard safety invariants.

\section{Benchmark Dataset and Ground Truth}
We designed a prototype benchmark spanning 9 destinations across two nature tourism categories in India: \textbf{Waterfalls} and \textbf{Trekking}. Provenance is incomplete and detailed in the accompanying data audit.

\subsection{Destination Selection}
The nine destinations in the project dataset span two activity categories and have supplied popularity values; those values have not yet been independently audited against external records:
\begin{itemize}
    \item \textbf{Waterfalls (5)}: Dudhsagar Falls (Goa), Jog Falls (Karnataka), Jogini Waterfall (Himachal Pradesh), Nohkalikai Falls (Meghalaya), Hogenakkal Falls (Tamil Nadu).
    \item \textbf{Trekking (4)}: Kedarkantha Trek (Uttarakhand), Valley of Flowers (Uttarakhand), Kashmir Great Lakes (KGL) Trek (Jammu and Kashmir), Chembra Peak (Kerala).
\end{itemize}

\subsection{Ground Truth Relevance Rubric}
Provisional relevance labels $g(d, m) \in \{0, 1, 2, 3\}$ were assigned by one project annotator using a written rubric. The labels have not been independently reviewed, and evidence supporting every destination-month rationale is not present in the repository:
\begin{itemize}
    \item \textbf{Grade 3 (Optimal)}: Peak seasonal condition marked open in the provisional inputs, with optimal flow/trail status and pleasant climate.
    \item \textbf{Grade 2 (Good)}: Safe, enjoyable visitation with minor compromises (e.g., slight pre-monsoon flow reduction, light rain requiring poncho).
    \item \textbf{Grade 1 (Marginal / Poor)}: Substantially degraded experience (e.g., waterfall reduced to a thin dry trickle on bare rock, dense fog with 0\% visibility, night camping banned).
    \item \textbf{Grade 0 (Unsuitable / Hazardous / Closed)}: Official administrative ban, life-safety hazard, avalanche risk, or lethal flash floods.
\end{itemize}

\section{Methodology}
We formulate three recommendation methods over candidate set $\mathcal{D}_a$ for activity category $a$ and query month $m$.

\subsection{Method A: Popularity-Based Baseline}
Method A mirrors commercial aggregators by ranking solely on historical normalized popularity $Pop(d) \in [0, 1]$:
\begin{equation}
S_A(d) = Pop(d)
\end{equation}
The ranking is completely time-invariant across all 12 months.

\subsection{Method B: Seasonal Multi-Factor Scoring}
Method B calculates an environmental fitness score combining weather comfort, scenic/flow viability, and popularity:
\begin{equation}
S_B(d, m, a) = w_1 S_{\text{scenic}}(d, m) + w_2 W(d, m, a) + w_3 Pop(d)
\end{equation}
where $W=0.5 TempFit+0.5 RainFit$ for trekking and $W=0.4 TempFit+0.6 RainFit$ for waterfalls. Temperature and rainfall fit use the piecewise thresholds implemented in the engine. Seasonal weights are $w_1=0.50$, $w_2=0.30$, and $w_3=0.20$.

\subsection{Method C: Disruption-Aware Seasonal Ranking}
Method C starts with the Method B seasonal score and applies only disruption records that pass evidence, event-date, destination, month, activity, and scope checks. A qualifying full-destination closure is removed from the actionable ranking and placed in a separate quarantine list. Applicable partial restrictions multiply the seasonal score by configured penalty factors. If no record qualifies, the Method C score equals the Method B score. This policy is an algorithmic property conditional on supplied data; it does not establish real-world safety.
\section{Empirical Evaluation and Results}
The regenerated Node benchmark contains 36 scenarios: 12 waterfall, 12 trekking, and 12 mixed-pool month scenarios. Requested K is 3 for single-activity scenarios and 5 for mixed-pool scenarios. Precision@K uses effective returned K. NDCG uses a shared ideal ranking over the original scenario candidate pool; relevance metrics are undefined if any candidate label is missing or IDCG is zero. Each method's eligible pool is recorded separately. Closure conflict is evaluated only against current evidence-backed full-closure records and is N/A when none qualifies.

\begin{table*}[t]
\centering
\caption{Exploratory recommendation results on provisional labels; K varies by scenario pool}
\label{tab:recsys_results}
\begin{tabular}{l|ccc|ccc|ccc}
\hline
\textbf{Metric} & \multicolumn{3}{c|}{\textbf{Category-specific (24)}} & \multicolumn{3}{c|}{\textbf{Mixed pool (12)}} & \multicolumn{3}{c}{\textbf{Overall (36)}} \\
& A & B & C & A & B & C & A & B & C \\
\hline
\textbf{Precision@K} & 0.6945 & 0.8611 & 0.8611 & 0.6667 & 0.9333 & 0.9333 & 0.6852 & 0.8852 & 0.8852 \\
\textbf{NDCG@K} & 0.6483 & 0.9789 & 0.9789 & 0.5462 & 0.9749 & 0.9749 & 0.6143 & 0.9776 & 0.9776 \\
\textbf{Full-closure conflict} & N/A & N/A & N/A & N/A & N/A & N/A & N/A & N/A & N/A \\
\hline
\end{tabular}
\end{table*}

The month split is exploratory because its selection rationale and independence from prior tuning are undocumented. Method B exceeds Method A on provisional labels. Method C ties Method B because zero of 11 supplied disruption records is evidence-backed; the active disruption path is not evaluated by this run, and no closure-conflict reduction can be inferred.

\section{Discussion and Limitations}
These results are descriptive and conditional on supplied values and single-annotator provisional labels. Because the rubric uses seasonal conditions also consumed by Methods B and C, the evaluation may reward agreement with the rubric rather than user preference or independently observed outcomes. Nine destinations are insufficient to establish generalization. Record-level evidence for monthly environmental values, popularity statistics, scenic indices, and disruptions is incomplete. The 108 relevance labels have no independent review or per-label supporting evidence. Seasonal weights and penalties are design choices, not empirically optimized. The holdout split is exploratory rather than confirmatory.

\section{Conclusion}
Across 36 scenarios, overall macro-average NDCG@K is 0.6143 for Method A and 0.9776 for Methods B and C. No closure-conflict rate is estimable because no supplied record passes the evidence policy. This benchmark does not demonstrate real-world safety or disruption-aware effectiveness. Further claims require a completed source audit, independent relevance judgments with adjudication, and evaluation under a predeclared protocol using evidence-backed disruption records.
\begin{thebibliography}{00}
\bibitem{gavalas2014} D. Gavalas, C. Konstantopoulos, K. Mastakas, and G. Pantziou, ``Mobile recommender systems in tourism,'' \textit{J. Netw. Comput. Appl.}, vol. 39, pp. 319--333, 2014.
\bibitem{werthner2004} H. Werthner and F. Ricci, ``E-commerce and tourism,'' \textit{Commun. ACM}, vol. 47, no. 12, pp. 101--105, 2004.
\bibitem{koren2009} Y. Koren, R. Bell, and C. Volinsky, ``Matrix factorization techniques for recommender systems,'' \textit{Computer}, vol. 42, no. 8, pp. 30--37, 2009.
\bibitem{abdollahpouri2017} H. Abdollahpouri, R. Burke, and B. Mobasher, ``Controlling popularity bias in learning-to-rank recommendation,'' in \textit{Proc. 11th ACM Conf. Recommender Syst.}, 2017, pp. 42--46.
\bibitem{adomavicius2011} G. Adomavicius and A. Tuzhilin, ``Context-aware recommender systems,'' in \textit{Recommender Systems Handbook}, Springer, 2011, pp. 217--253.
\bibitem{baltrunas2012} L. Baltrunas, B. Ludwig, S. Peer, and F. Ricci, ``Context relevance assessment and exploitation in mobile recommender systems,'' \textit{Pers. Ubiquitous Comput.}, vol. 16, no. 5, pp. 507--526, 2012.
\end{thebibliography}

\end{document}
